% Part: model-theory
% Chapter: interpolation
% Section: introduction

\documentclass[../../../include/open-logic-section]{subfiles}

\begin{document}

\olfileid[hi]{mod}{int}{int}

\olsection{परिचय}

अंतर्वेशन प्रमेय निम्न परिणाम है: मान लें कि
$\Entails !A \lif !B$। तब ऐसा !!a{sentence} $!C$ है कि
$\Entails !A \lif !C$ और $\Entails !C \lif !B$। इसके अतिरिक्त, प्रत्येक
!!{constant}, !!{function} और !!{predicate} ($\eq$ के अतिरिक्त), जो
$!C$ में आता है, $!A$ और~$!B$ दोनों में आता है। !!{sentence} $!C$
को $!A$ और~$!B$ का \emph{अंतर्वेशक} कहते हैं।

अंतर्वेशन प्रमेय अपने आप में रोचक है, किंतु उसका मुख्य महत्त्व इस बात में
है कि उसका उपयोग किसी सिद्धांत में परिभाषणीयता के बारे में परिणाम सिद्ध
करने और उन शर्तों को निर्धारित करने में किया जा सकता है जिनके अंतर्गत दो
संगत सिद्धांतों को मिलाने पर संगत सिद्धांत मिलता है। पहला परिणाम बेथ के
परिभाषणीयता प्रमेय के नाम से और दूसरा रॉबिन्सन के संयुक्त संगतता प्रमेय के
नाम से जाना जाता है।

\end{document}
